% Exact example: R=2; P=(4,0); T=(1,+/-sqrt(3)); polar x=1.
\begin{tikzpicture}[x=1.12cm,y=1.12cm,font=\small]
\draw[gray!45] (-2.3,0)--(4.25,0);
\draw[gray!75,thick] (0,0) circle (2);
\draw[orange!85!black,thick] (1,{sqrt(3)})--(4,0)--(1,{-sqrt(3)});
\draw[blue!65!black,very thick] (1,-2.25)--(1,2.3);
\fill (0,0) circle (1.4pt) node[below left] {$O$};
\fill[orange!85!black] (4,0) circle (1.7pt) node[right] {$P(4,0)$};
\fill (1,{sqrt(3)}) circle (1.7pt) node[above left,yshift=1mm] {$T_1$};
\fill (1,{-sqrt(3)}) circle (1.7pt) node[below right,xshift=1mm] {$T_2$};
\node[blue!65!black,above] at (1,2.35) {$l_P:x=1$};
\node[orange!85!black,above,rotate=-30] at (2.95,.85) {切线};
\node at (-1.1,1) {$R=2$};
\draw[CStatement!35,rounded corners=2mm] ([xshift=-3.5mm,yshift=-3.5mm]current bounding box.south west) rectangle ([xshift=3.5mm,yshift=3.5mm]current bounding box.north east);
\end{tikzpicture}
